% === Section 10: Cross-Domain Synthesis ===
\section{Cross-Domain Synthesis}\label{sec:crossdomain}

The active ORIUS defense lane is a literal three-domain program:
\textbf{Battery} as the witness row, plus bounded defended rows for
\textbf{Autonomous Vehicles} and \textbf{Healthcare Monitoring}.  The role of this
section is not to assert equal-depth theorem closure across all three rows.  It is
to show that the same \dcs{} kernel, contract stack, and governance surface survive
domain change while the evidence tier of each row is stated explicitly.

\subsection{Domain Roles}

\begin{table}[h]
\caption{Active three-domain roles and evidence posture.}
\label{tab:cross_domain_roles}
\centering\small
\begin{tabular}{llll}
\toprule
\textbf{Domain} & \textbf{Role} & \textbf{Runtime denominator} & \textbf{Bounded claim} \\
\midrule
Battery & witness row & full CPSBench witness episode & deepest theorem-to-artifact closure \\
AV & defended bounded row & 9{,}348 promoted runtime steps & TTC + predictive-entry-barrier contract \\
Healthcare & defended bounded row & bounded healthcare evidence lane & monitoring / alert-release semantics \\
\bottomrule
\end{tabular}
\end{table}

\subsection{Metric Comparison}

\Cref{tab:cross_domain_metrics_three} summarises the promoted metrics that are
shared across the three active rows.

\begin{table}[h]
\caption{Promoted three-domain metric comparison.}
\label{tab:cross_domain_metrics_three}
\centering\small
\begin{tabular}{lccc}
\toprule
\textbf{Metric} & \textbf{Battery} & \textbf{AV} & \textbf{Healthcare} \\
\midrule
Baseline TSVR           & 0.0393 & 0.1250 & 0.2917 \\
ORIUS TSVR              & 0.0000 & 0.0417 & 0.0417 \\
Absolute TSVR delta     & 0.0393 & 0.0833 & 0.2500 \\
Intervention rate       & 0.0208 & 0.9018 & 0.0542 \\
Certificate-valid release rate & 0.9931 & 1.0000 & 1.0000$^\dagger$ \\
Grouped coverage (low / mid / high) & 0.938 / 1.000 / 0.906 & 0.994 / 0.997 / 0.999 & 0.611 / 0.953 / 0.844 \\
P95 runtime latency (ms) & 0.812 & 0.744 & 0.701 \\
\bottomrule
\end{tabular}
\end{table}

{\footnotesize $^\dagger$ Healthcare currently reports certificate-valid release
via the governance-pass proxy on the promoted runtime export.  The paper keeps that
semantics explicit and does not overstate it as a direct CVA measurement.}

\subsection{Analysis}

\paragraph{Battery witness row.}
Battery remains the deepest reference row.  It is the only domain for which the
theorem narrative, runtime witness, and locked artifact chain are intended to carry
the full constructive story end to end.  The promoted battery row achieves
$\TSVR = 0$, but the purpose of the witness row is not leaderboard dominance; it is
the cleanest theorem-to-code-to-artifact closure for ORIUS under degraded
observation.

\paragraph{Autonomous Vehicles bounded row.}
AV is promoted under the longitudinal TTC + predictive-entry-barrier contract only.
ORIUS reduces the bounded validation-harness TSVR from 0.1250 to 0.0417, while the
full-corpus AV runtime summary remains locked separately.  This row demonstrates
that the same runtime contract can survive domain transfer to a faster,
interaction-heavy prediction setting, but it does \emph{not} claim full autonomous
driving closure.

\paragraph{Healthcare bounded row.}
Healthcare is promoted only as a bounded monitoring and alert-release row.  ORIUS
reduces the bounded TSVR from 0.2917 to 0.0417, with grouped calibration and
governance evidence reported explicitly.  The grouped coverage remains weakest in
the low-reliability bucket, which is exactly why the prose keeps the healthcare row
bounded and avoids any clinical-deployment claim.

\subsection{Kernel Invariance}

The cross-domain claim that ORIUS \emph{does} defend is kernel invariance under
typed adaptation:
\begin{enumerate}[nosep]
  \item The same \texttt{dc3s} runtime kernel is used across Battery, AV, and
        Healthcare.
  \item The same contract stack and governance surface remain active across all
        three domains.
  \item Domain-specific behavior is isolated in the adapter layer and the
        domain-local safety contract.
\end{enumerate}

This is the correct level of universality for the active branch: \emph{runtime and
governance universality}, not equal-depth empirical closure across every possible
physical-AI domain.

\subsection{What Transfers and What Does Not}

\paragraph{What transfers.}
The five-stage \dcs{} pipeline transfers cleanly: degradation detection, interval
calibration, action-set tightening, repair/fallback, and certificate issuance all
appear on each active row.

\paragraph{What does not transfer automatically.}
The witness strength does not transfer automatically.  Battery remains deeper than
AV and Healthcare.  The defended theorem core therefore treats the rows
asymmetric\-ally: Battery is the constructive witness, while AV and Healthcare are
bounded defended rows that validate the runtime contract under different observation
and actuation semantics.

\subsection{Implications}

The three-domain synthesis supports the flagship ORIUS claim in its reviewer-safe
form: OASG is a runtime hazard of degraded observation, and ORIUS provides a
reliability-aware safety layer that can be operationalized across Battery, AV, and
Healthcare without changing the kernel semantics.  The synthesis does \emph{not}
claim equal-depth theorem closure, a new universal controller, or unrestricted
deployment readiness in any of the bounded rows.
